% CLAUDE REPORT::BRIDGE-0001::2026-11-12_1049_PST
% Shared worst-case derivation objective: what the conditional bridge
%   SAT in P  ==>  \overline W_R(N) <= poly(N)
% can and cannot be, for a fixed refutation calculus R.
% UTF-8, XeLaTeX/pdfLaTeX compatible.
\documentclass[11pt]{article}
\usepackage[margin=1in]{geometry}
\usepackage{amsmath,amssymb,amsthm}
\usepackage{hyperref}
\usepackage{url}
\newtheorem{theorem}{Theorem}
\newtheorem{lemma}[theorem]{Lemma}
\newtheorem{corollary}[theorem]{Corollary}
\newtheorem{proposition}[theorem]{Proposition}
\theoremstyle{definition}
\newtheorem{definition}[theorem]{Definition}
\newcommand{\Pclass}{\mathsf{P}}
\newcommand{\NP}{\mathsf{NP}}
\newcommand{\coNP}{\mathsf{coNP}}
\newcommand{\PSPACE}{\mathsf{PSPACE}}
\newcommand{\SAT}{\mathrm{SAT}}
\newcommand{\UNSAT}{\mathrm{UNSAT}}
\newcommand{\Wb}{\overline W}
\newcommand{\PB}{\mathrm{PB}}
\title{CLAUDE BRIDGE-0001: The Conditional Bridge for a Fixed Calculus --- Truth Table, Monotonicity, and a Relativization Barrier}
\author{P vs NP Research Project --- Claude scheduled research run}
\date{12 November 2026, 10:49 PST}
\begin{document}
\maketitle

\section*{Status}
\textbf{P versus NP remains OPEN.} Nothing below proves or refutes it.
Lemma~\ref{lem:truth}, Theorem~\ref{thm:conservation}, Theorem~\ref{thm:mono},
Theorem~\ref{thm:relativ} and Lemma~\ref{lem:auto} are \textbf{PROVED/EXACT}
(elementary). Corollary~\ref{cor:known} and \S\ref{sec:ef} combine these with
\textbf{KNOWN THEOREMS} from the literature. The bridge $B(R)$ for every
calculus whose worst-case minimum is not yet known to be superpolynomial
(Extended Frege, Frege, the project calculus $R_{\mathrm{mix341}}$) is
\textbf{OPEN}. No finite experiment is claimed; the only executed code is the
ledger regeneration in \S\ref{sec:ledger}.

\section*{Target (shared five-role objective)}
For a fixed Cook--Reckhow refutation system $R$ with cost $\operatorname{cost}_R$
polynomially related to proof size,
\[
\Wb_R(N)=\max_{\substack{F\in\UNSAT\\|\langle F\rangle|\le N}}\
\min_{\pi:F\vdash_R\bot}\operatorname{cost}_R(\pi),
\qquad
\PB(R):\ \exists c\ \forall N\ \Wb_R(N)\le N^c+c,
\]
\[
B(R):\qquad \SAT\in\Pclass\ \Longrightarrow\ \PB(R).
\]
The plan is: prove $\neg\PB(R)$ and $B(R)$ for the same fixed $R$, which gives
$\SAT\notin\Pclass$. The optimal derivation is granted for free; this is a
valid lower-bound relaxation for any procedure that executes an $R$-derivation.
This run asks what $B(R)$ is logically, how it moves with $R$, and which proof
techniques cannot establish it.

\section{Barrier-transfer literature preflight}
\textbf{Main obstacles.} (O1) $B(R)$ asks for an \emph{oracle-free} conclusion
(short proofs in a fixed calculus) from a hypothesis about an arbitrary
algorithm. (O2) Every calculus with a proved superpolynomial lower bound is weak
enough that some polynomial-time algorithm defeats its hard family.

\textbf{Sources checked this run.}
\begin{itemize}
\item Cook--Reckhow 1979: $\NP=\coNP$ iff some propositional proof system is
polynomially bounded.
\item Baker--Gill--Solovay 1975: $\Pclass^O=\NP^O$ for any $\PSPACE$-complete $O$.
\item Aaronson--Wigderson 2009 (algebrization): there are $A,\tilde A$ with
$\NP^{\tilde A}\subseteq\Pclass^{A}$ ($A$ $\PSPACE$-complete, $\tilde A$ a
low-degree extension).
\item Cook 1975 and Krajíček--Pudlák 1990: $\forall\Pi^b_1$ theorems of $PV$
translate to polynomial-size Extended Frege proofs.
\item Atserias--Müller, JACM 2020: Resolution is not automatable unless
$\Pclass=\NP$.
\item Pich--Santhanam, \emph{Towards $\Pclass\ne\NP$ from Extended Frege lower
bounds} (ECCC TR23-199, arXiv:2312.08163; journal version cited in the corpus
as JACM 2026). Under self-provability or bounded-arithmetic hypotheses,
superpolynomial lower bounds for EF (or EF plus extra axioms) imply
$\Pclass\ne\NP$ or $\NP\not\subseteq\Pclass/\mathrm{poly}$.
\item Recent relativized optimal-proof-system work (Egidy--Glaßer and coauthors,
arXiv:2408.07408 and arXiv:2602.02294) builds oracles that separate questions
about optimal proof systems. That confirms relativized proof-system statements
are delicate. It does not supply a fixed-$R$ bridge.
\end{itemize}
Titles and years are recalled from memory except where a URL was fetched
(arXiv:2312.08163 abstract was read this run). The technique tags in the
ledger stay \texttt{n.a.}/\texttt{Unverified} accordingly.

\textbf{Transfer classification.} The oracle-world construction (complexity
theory $\to$ proof complexity) is an \emph{exact object transfer}: the
proof system is fixed, so its proof lengths do not depend on the oracle
(Theorem~\ref{thm:relativ}). The worst-case/average-case architecture of LWE
(PQ ledger item 5, ``extract the resource from black-box solver behaviour'')
is a \emph{failed transfer} for every calculus in Corollary~\ref{cor:known}.
Any black-box extraction argument relativizes, and Theorem~\ref{thm:relativ}
rules relativizing arguments out. Schur-square, collision-entropy and hash
ideas (items 1--4, 6--7) do not bear on O1 or O2; no transfer is claimed.

\section{The bridge is a disjunction}
\begin{lemma}[Truth table of the bridge]\label{lem:truth}
For every fixed $R$: $B(R)\iff(\SAT\notin\Pclass\ \vee\ \PB(R))$.
\end{lemma}
\begin{proof}
$X\Rightarrow Y$ is equivalent to $\neg X\vee Y$.
\end{proof}

\begin{theorem}[Conservation: the decomposition is never easier than the goal]\label{thm:conservation}
For every fixed $R$,
\[
B(R)\wedge\neg\PB(R)\iff \SAT\notin\Pclass\wedge\neg\PB(R).
\]
In particular, once $\neg\PB(R)$ is proved, $B(R)$ is equivalent to
$\Pclass\neq\NP$.
\end{theorem}
\begin{proof}
Use Lemma~\ref{lem:truth}. Under $\neg\PB(R)$ the disjunction reduces to
$\SAT\notin\Pclass$.
\end{proof}

So the two-step plan always costs \emph{at least} a proof of
$\Pclass\neq\NP$, plus the extra obligation $\neg\PB(R)$. Its value can only be
methodological: a structured route into the separation, not a weaker statement.
Codex 2026-11-11 1203 §``A weaker sufficient bridge'' observed the
quasipolynomial form of this; Theorem~\ref{thm:conservation} is the polynomial
form, adopted here as the baseline.

\begin{corollary}[Known-hard calculi]\label{cor:known}
Each of the following calculi has an unconditional superpolynomial lower
bound on unsatisfiable CNF families, so for each of them $B(R)$ \emph{is}
$\Pclass\neq\NP$: tree-like and general Resolution (Haken 1985); Cutting
Planes (Pudlák 1997); Polynomial Calculus over any field (Razborov 1998;
Impagliazzo--Pudlák--Sgall 1999); depth-$d$ Frege for each fixed $d$ (Ajtai
1988; Pitassi--Beame--Impagliazzo 1993; Krajíček--Pudlák--Woods 1995).
\end{corollary}
\begin{proof}
These are KNOWN THEOREMS giving $\neg\PB(R)$. Apply Theorem~\ref{thm:conservation}.
\end{proof}

\section{Monotonicity: bridge and lower bound move in opposite directions}
Write $R\le_p R'$ if $R'$ p-simulates $R$ in proof size.
\begin{theorem}[Opposite monotonicity]\label{thm:mono}
If $R\le_p R'$ then $B(R)\Rightarrow B(R')$ and $\neg\PB(R')\Rightarrow\neg\PB(R)$.
If $R_{\mathrm{mix}}$ contains $R_1,\dots,R_k$ as subsystems, each step charged at
most a fixed polynomial of its native cost, then for every UNSAT $F$
\[
\min_{\pi:F\vdash_{R_{\mathrm{mix}}}\bot}\operatorname{cost}(\pi)
\le q\Bigl(\min_i\ \min_{\pi:F\vdash_{R_i}\bot}\operatorname{cost}_{R_i}(\pi)\Bigr)
\]
for a fixed polynomial $q$.
\end{theorem}
\begin{proof}
$\PB(R)$ implies $\PB(R')$ by composing the simulation polynomial with the
bound. Both implications follow, the second by contraposition. For the mixed
system, take the cheapest component refutation and embed it.
\end{proof}

\textit{Reading.} Strengthening the calculus makes the bridge easier and the
lower bound harder. Adaptive strategy switching (priority 4 of the directive)
can only \emph{lower} the minimum: a hard instance must beat every component on
the \emph{same} source. The corpus Tseitin/XOR control (LSPV-1/334) and the
PHP/Cutting-Planes control are instances of this.

\begin{proposition}[Endpoint: optimal calculus; ADOPTED from LSPV-1/336]
If $R^\ast$ is size-optimal (p-simulates every Cook--Reckhow system in length),
then $B(R^\ast)$ holds, and $\neg\PB(R^\ast)\iff\NP\neq\coNP$.
\end{proposition}
\begin{proof}
$\SAT\in\Pclass$ gives $\NP=\coNP$. Cook--Reckhow then gives a polynomially
bounded system, which $R^\ast$ simulates. Conversely, $\neg\PB(R^\ast)$ means no
system is polynomially bounded, by optimality.
\end{proof}

So the spectrum runs from weak calculi, where the lower bound is known and the
bridge is all of $\Pclass\ne\NP$, to an optimal calculus, where the bridge is
free and the lower bound is $\NP\ne\coNP$, strictly at least as hard. Both
obligations are open only at intermediate strength (Frege, EF, and conditionally
$R_{\mathrm{mix341}}$ by LSPV-1/592).

\section{A relativization and algebrization barrier for $B(R)$}\label{sec:relativ}
For an oracle $O$ keep $R$ \emph{fixed and oracle-free}, exactly as the
directive does (one preselected calculus, fixed weights). Define
\[
B^O(R):\qquad \NP^O\subseteq\Pclass^O\ \Longrightarrow\ \PB(R).
\]
Call a proof of $B(R)$ \emph{relativizing} if the same argument proves
$B^O(R)$ for every $O$.

\begin{theorem}[Relativization barrier]\label{thm:relativ}
Let $O$ be $\PSPACE$-complete. Then $B^O(R)\iff\PB(R)$. Consequently:
\begin{enumerate}
\item a relativizing proof of $B(R)$ proves $\PB(R)$, hence $\NP=\coNP$;
\item for every $R$ in Corollary~\ref{cor:known}, $B(R)$ has no relativizing
proof (unconditionally);
\item for every $R$, if $\NP\neq\coNP$, then $B(R)$ has no relativizing proof.
\end{enumerate}
The same holds with algebrizing proofs: for the Aaronson--Wigderson pair
$(A,\tilde A)$, the algebrized bridge
$\NP^{\tilde A}\subseteq\Pclass^{A}\Rightarrow\PB(R)$ is again equivalent to
$\PB(R)$.
\end{theorem}
\begin{proof}
$\PB(R)$ makes no reference to $O$: proofs, verifier and $\Wb_R$ are
oracle-free. By Baker--Gill--Solovay the antecedent $\NP^O\subseteq\Pclass^O$
holds, so $B^O(R)$ reduces to its consequent. Item 1: a relativizing argument
yields $B^O(R)$ for this $O$, hence $\PB(R)$; Cook--Reckhow gives $\NP=\coNP$.
Item 2: $\PB(R)$ is false there. Item 3: $\PB(R)$ implies $\NP=\coNP$. The
algebrization case is identical, using the known pair with
$\NP^{\tilde A}\subseteq\Pclass^A$.
\end{proof}

\textit{Convention firewall.} If instead the calculus is relativized too
($R^O$ may cite oracle answers as axioms), then $R^O$ is trivially polynomially
bounded for $\PSPACE$-complete $O$ and $B^O(R^O)$ holds vacuously. The barrier
is about the project's convention, a fixed oracle-free calculus. It
says the argument must use the \emph{code} or resource bound of the hypothetical
algorithm in a way that does not survive adding a $\PSPACE$ oracle. Black-box
use of the decider's answers, query counts, or input/output behaviour is not
enough.

\textit{Consequence for the five roles.} Any proposed transfer theorem with
the shape ``given any SAT decider (used only through its answers or its
runtime), construct an $R$-derivation of comparable cost'' relativizes. For the
calculi in Corollary~\ref{cor:known} it is therefore false outright. This
generalizes the corpus counterexamples to $\mathrm{TR}(\mathrm{Res})$
(LSPV-1/334, Tseitin versus Gaussian elimination) from one family to all
relativizing arguments.

\section{Free optimal sequence: exactly realized under the hypothesis}
\begin{lemma}[Automation under $\SAT\in\Pclass$]\label{lem:auto}
If $\SAT\in\Pclass$, then for every Cook--Reckhow $R$ there is a polynomial-time
algorithm that, given $F$ and $1^S$, outputs an $R$-refutation of $F$ of size
at most $S$ whenever one exists. Hence under the hypothesis some $R$-executing
SAT procedure runs in time $\mathrm{poly}(N+\min_\pi|\pi|)$ on every UNSAT input.
\end{lemma}
\begin{proof}
The language $\{(F,1^S,p):p\text{ extends to an }R\text{-refutation of size}\le S\}$
is in $\NP$, hence in $\Pclass$. Build the refutation bit by bit
(search-to-decision).
\end{proof}

So the directive's principle --- not charging for discovery of the optimal
sequence --- costs nothing in the world where $\SAT\in\Pclass$. Then
$\Wb_R$ \emph{is} the running time of the best $R$-based solver, up to a
polynomial, and $B(R)$ says exactly that this solver is polynomial. In
$\Pclass\ne\NP$ worlds, discovery is genuinely expensive for Resolution
(Atserias--Müller). Not charging it there is a real relaxation, but a legitimate one
for lower bounds, as the directive states.

\section{The only known code-using bridge (Extended Frege)}\label{sec:ef}
\textbf{KNOWN THEOREM} (Cook 1975; Krajíček--Pudlák 1990). Suppose $A$ runs in
polynomial time and $PV\vdash\forall F\,\forall\alpha\,(A(F)=0\rightarrow
\neg\mathrm{Sat}(F,\alpha))$. Then every UNSAT $F$ has an EF refutation of size
$\mathrm{poly}(|F|)$. Its two parts are the translation of the $PV$ proof and
an EF verification of the concrete run $A(F)=0$. Hence
\[
B_{PV}(\mathrm{EF}):\quad
\exists A\ \bigl(A\text{ poly-time, } PV\vdash\mathrm{Sound}_A,\ A\text{ complete}\bigr)
\Longrightarrow\PB(\mathrm{EF})
\]
is a theorem. The argument uses the code of $A$, consistent with
Theorem~\ref{thm:relativ}. The gap to $B(\mathrm{EF})$ is exactly that
$\SAT\in\Pclass$ does not supply $PV$-provable soundness. Pich--Santhanam work
in this gap. By Theorem~\ref{thm:mono}, any calculus that p-simulates ER
inherits $B_{PV}$. LSPV-1/592 shows $R_{\mathrm{mix341}}$ would p-simulate ER,
conditional on its unauthenticated native rule interface, so the same holds
for $R_{\mathrm{mix341}}$. That route also makes the needed lower bound at
least an EF lower bound, which is open.

\textbf{Recommended target (CANDIDATE).} Replace the open $B(R)$ by the
already-proved $B_{PV}(\mathrm{EF})$ plus a theory-level obligation:
\[
\text{(T)}\qquad \SAT\in\Pclass\ \Longrightarrow\ \exists\text{ poly-time }A
\text{ correct for SAT with } PV\vdash\mathrm{Sound}_A .
\]
Then (T) and $\neg\PB(\mathrm{EF})$ together give $\Pclass\neq\NP$. (T) is a
statement about \emph{provability}, not proof length. It is not refuted by
Theorem~\ref{thm:relativ}, since it is not oracle-free in the relevant sense.
(T) is \textbf{OPEN}, and Theorem~\ref{thm:conservation} still applies: (T) is
at least as hard as the separation once $\neg\PB(\mathrm{EF})$ is known.

\section{Hostile controls and red team}
\begin{itemize}
\item Tseitin on expanders: Resolution exponential, Gaussian elimination
polynomial. This is consistent with Corollary~\ref{cor:known}: $B(\mathrm{Res})$ is not
refuted by it (it is $\Pclass\ne\NP$), but every \emph{pointwise}
algorithm-to-derivation simulation into Resolution is (LSPV-1/334, retained).
\item PHP: Resolution exponential, Cutting Planes polynomial. This instantiates
Theorem~\ref{thm:mono} under strategy switching.
\item Relativized-calculus convention: $B^O(R^O)$ is vacuous. Recorded as a
firewall, not a counterexample.
\item Machine-trace soundness (LSPV-1/334): a trace of $A$ ending in
``UNSAT'' is not a proof unless the soundness of $A$ is itself justified. This is
precisely the $PV\vdash\mathrm{Sound}_A$ premise of \S\ref{sec:ef}.
\item Natural proofs: no circuit lower bound or constructive property is
claimed. Not invoked.
\end{itemize}

\section*{CORPUS REVIEWED}
Root \path{AGENTS.md} (including the 2026-11-09 knot/link handoff note and the
world-ledger pointer); recent root \path{main.tex} ownership markers;
\path{research/agents/codex/2026-11-11_1203_SAT_POLYNOMIAL_ASSUMPTION_AUDIT.tex}
(bridges (B), (12)); \path{research/agents/codex/2026-11-11_1331_TRANSFER_RECOVERY_AND_RATE_AMPLIFICATION.tex}
(relativization paragraph);
\path{research/agents/lspv-1/2026-11-12_0703_LSPV1_ER_STRENGTH_BARRIER592.md};
\path{research/agents/lspv-1/2026-11-02_1100_LSPV-1_FREE_CHOICE_TRANSFER_BARRIER334.md};
\path{research/agents/lspv-1/2026-11-02_1256_LSPV-1_CLOCKED_OPTIMALITY_REFLECTION336.tex};
the coordination handoff list through LSPV-4/566 (2026-11-12 10:31);
the Claude folder and world ledger. Repository searches covered
``optimal proof system'', ``automatab'', ``relativiz'', ``bridge'' and
$\Wb_R$. This is a targeted survey, not a byte-level reread of all
historical files.

\section*{DUPLICATION / CONTRADICTION AUDIT}
\begin{itemize}
\item Theorem~\ref{thm:conservation}: the quasipolynomial form is already in
Codex 1203. It is \emph{duplicated} here in polynomial form and marked as such.
\item Optimal-calculus endpoint: already LSPV-1/336. ADOPTED with attribution.
\item New this run: Theorem~\ref{thm:relativ} (fixed-calculus relativization and
algebrization barrier). Codex 1331's relativization note concerns a different,
oracle-matched computation bridge, and says only that its transfers do not
bypass relativization. Also new: Theorem~\ref{thm:mono} (opposite monotonicity),
Corollary~\ref{cor:known} as an explicit list, Lemma~\ref{lem:auto} tied to the
directive's free-optimal-sequence principle, and target (T).
\item No contradiction found with the root \path{main.tex} status (``P versus
NP remains OPEN'', B1/B2/B3 OPEN).
\end{itemize}

\section*{WHAT CHANGED RELATIVE TO ROOT MAIN}
Root \path{main.tex} is \emph{not edited} by this run. It has ownership regions
only for LSPV-1..5, and this Claude run owns none of them. Its OPEN status is
unchanged. Suggested input for the LSPV-5 arbiter: record that $B(R)$ for any
fixed oracle-free $R$ cannot have a relativizing or algebrizing proof unless
$R$ is polynomially bounded. Record also that for every currently lower-bounded
calculus, $B(R)$ is literally $\Pclass\ne\NP$.

\section*{ADOPTED / REJECTED / RETIRED PRIOR COMPONENTS}
ADOPTED: LSPV-1/334 counterexample to pointwise $\mathrm{TR}(\mathrm{Res})$;
LSPV-1/336 optimal-calculus bridge; LSPV-1/592 conditional ER simulation;
Codex 1203 quasipolynomial bridge.
REJECTED (as a route to $\Pclass\ne\NP$ for weak calculi): any black-box
solver-to-derivation extraction (Theorem~\ref{thm:relativ}).
RETIRED: none.

\section*{STALE OR OVERSTATED CLAIMS FOUND}
None found in the files reviewed. One caution for future text: describing
$B(R)$ as ``the weakest sufficient bridge'' is accurate as a sufficient
condition, but by Theorem~\ref{thm:conservation} it is never logically weaker
than $\Pclass\neq\NP$ once the lower bound for $R$ is in hand.

\section*{Predecessor write check}
This Claude run is outside the LSPV-1..5 stagger and has no cyclic predecessor.
The most recent LSPV-4/566 commits on \texttt{main} (2026-11-12 10:31 handoff)
were present on \texttt{origin/main} at fetch time.

\section{Ledger update}\label{sec:ledger}
Rows L41--L43 were added to
\path{research/agents/claude/world-ledger/build_ledger.py}, and the CSV/XLSX
were regenerated (43 rows; one run under 30 s). Automatic categories:
L41 $B(\mathrm{Res})$ --- \emph{Equivalent to $\Pclass\ne\NP$};
L42 Resolution automatable --- \emph{Equivalent to $\Pclass=\NP$};
L43 $B(\mathrm{EF})$ --- \emph{Half-win (true if $\Pclass\ne\NP$)}.
All consistency checks are OK.

\section*{Next gap}
Either (a) prove a nonrelativizing, code-using bridge for a calculus whose
lower bound is open (the target is (T), or a Pich--Santhanam-style
self-provability statement), or (b) prove $\neg\PB(R)$ for a calculus at least as
strong as EF. Both are \textbf{OPEN}. \textbf{P versus NP remains OPEN.}

% END_OF_VERIFIED_REPORT::CLAUDE::BRIDGE-0001
\end{document}
